\documentclass[11pt]{article}

\usepackage{amsmath}
\usepackage{amssymb}
\usepackage{amsthm}
\usepackage{mathtools}
\usepackage{enumitem}
\usepackage[a4paper,margin=1in]{geometry}
\usepackage{xcolor}
\usepackage{hyperref}
\usepackage{longtable}
\usepackage{array}
\usepackage{booktabs}

\hypersetup{
    colorlinks=true,
    linkcolor=blue!50!black,
    urlcolor=blue!60!black,
    citecolor=blue!50!black
}

\title{\textbf{ISI MSQE Solutions}\\[4pt]
\large Paper: PEA \\
\large Year/Batch: 2026 \\
\large Prepared for: Statstrive}
\author{}
\date{}

\setlength{\parindent}{0pt}
\setlength{\parskip}{4pt}

\begin{document}

\maketitle
\thispagestyle{empty}

\vspace{0.6cm}

\begin{center}
\large \textbf{ISI MSQE 2026 --- PEA Question Paper: Polished Solutions}
\end{center}

\vspace{0.4cm}

This document presents complete, exam-quality solutions to all 30 multiple-choice
questions of the ISI MSQE 2026 PEA paper. Each solution is derived from first
principles using standard tools of calculus, real analysis, probability,
statistics, microeconomics and macroeconomics.

\newpage

\section*{Answer Key Summary (PEA)}

\begin{center}
\begin{tabular}{c l c c}
\toprule
Question & Topic & Correct Option & Status \\
\midrule
1  & Sequences and Series (Telescoping)         & A & Verified \\
2  & Algebra                                    & C & Verified \\
3  & Calculus (Definite Integrals)              & D & Verified \\
4  & Calculus (Definite Integrals)              & C & Verified \\
5  & Real Analysis (Rolle / MVT)                & C & Verified \\
6  & Real Analysis (Cauchy MVT)                 & B & Verified \\
7  & Real Analysis (Integrals)                  & C & Verified \\
8  & Calculus (L'H\^opital, Symmetry)           & B & Verified \\
9  & Statistics (Correlation)                   & B & Verified \\
10 & Regression                                 & D & Verified \\
11 & Hypothesis Testing                         & D & Verified \\
12 & Random Variables                           & B & Verified \\
13 & Statistics (Unbiased Estimation)           & D & Verified \\
14 & Statistics (Descriptive)                   & B & Verified \\
15 & Probability (Bayes)                        & B & Verified \\
16 & Consumer Theory                            & D & Verified \\
17 & General Equilibrium (Exchange)             & A & Draft \\
18 & General Equilibrium (Public Goods)         & -- & Needs Review \\
19 & Market Structure (Monopoly Pricing)        & B & Verified \\
20 & Game Theory (Public Good NE)               & B & Verified \\
21 & Public Goods (Samuelson)                   & C & Verified \\
22 & Tax Incidence                              & C & Verified \\
23 & Consumer Theory (Engel Aggregation)        & A & Verified \\
24 & Consumer Theory (Aggregate Elasticity)     & C & Verified \\
25 & International Trade (Cobb--Douglas)        & D & Verified \\
26 & Macroeconomics (Classical)                 & C & Verified \\
27 & Macroeconomics (Fixed Price)               & D & Verified \\
28 & Growth (Solow)                             & A & Verified \\
29 & Growth (Solow, Cross-country)              & B & Verified \\
30 & Intertemporal Choice                       & B & Verified \\
\bottomrule
\end{tabular}
\end{center}

\bigskip

\section*{Topic Summary}

\begin{center}
\begin{tabular}{l c}
\toprule
Topic Area & Questions \\
\midrule
Calculus / Integration                & 3, 4, 8 \\
Real Analysis (MVT, Rolle)            & 5, 6, 7 \\
Algebra / Sequences                   & 1, 2 \\
Probability \& Random Variables       & 12, 15 \\
Statistics (Estimation, Descriptive)  & 9, 13, 14 \\
Regression \& Hypothesis Testing      & 10, 11 \\
Consumer Theory                       & 16, 23, 24 \\
General Equilibrium \& Public Goods   & 17, 18, 21 \\
Market Structure                      & 19 \\
Game Theory                           & 20 \\
Tax Incidence                         & 22 \\
International Trade                   & 25 \\
Macroeconomics \& Money               & 26, 27 \\
Growth Theory                         & 28, 29 \\
Intertemporal Choice                  & 30 \\
\bottomrule
\end{tabular}
\end{center}

\newpage

\section*{Full Polished Solutions}

%--------------------------------------------------
\section*{Question 1}
\textbf{Paper:} PEA \quad \textbf{Year:} 2026 \quad \textbf{Topic:} Sequences and Series \\
\textbf{Subtopic/Concept:} Telescoping sum, partial fractions \quad
\textbf{Difficulty:} Easy \quad \textbf{Status:} Verified

\subsection*{Question}
Suppose $a$ and $b$ are positive integers that are relatively prime, and
\[
\frac{a}{b} \;=\; \sum_{n=1}^{2026} \frac{1}{n^2 + 15n + 56}.
\]
Then the value of $a+b$ is
\begin{enumerate}[label=(\Alph*)]
\item 9149 \item 9229 \item 8189 \item 11
\end{enumerate}

\subsection*{Solution}
Factor the denominator:
\[
n^2 + 15n + 56 \;=\; (n+7)(n+8).
\]
Partial fractions give
\[
\frac{1}{(n+7)(n+8)} \;=\; \frac{1}{n+7} - \frac{1}{n+8}.
\]
Hence the sum telescopes:
\[
\sum_{n=1}^{2026}\left(\frac{1}{n+7} - \frac{1}{n+8}\right)
= \frac{1}{8} - \frac{1}{2034}
= \frac{2034 - 8}{8 \cdot 2034}
= \frac{2026}{16272}.
\]
Now simplify $\dfrac{2026}{16272}$. Since $2026 = 2 \cdot 1013$ and $16272 = 2 \cdot 8136$,
\[
\frac{2026}{16272} = \frac{1013}{8136}.
\]
We verify $\gcd(1013, 8136)=1$: $8136 = 2^3 \cdot 3^2 \cdot 113$, and $1013$ is prime
(no prime factor up to $\lfloor\sqrt{1013}\rfloor = 31$ divides it). Thus the fraction
is already in lowest terms with $a = 1013$, $b = 8136$.

Therefore
\[
a + b = 1013 + 8136 = 9149.
\]

\subsection*{Final Answer}
\[
\boxed{\text{Option A: } 9149}
\]

\subsection*{Common Trap / Note}
A common error is to take $1/8 - 1/(N+8)$ with $N=2026$ but then forget to verify
coprimality. The fraction $\tfrac{2026}{16272}$ reduces by a factor of $2$ only.

%--------------------------------------------------
\section*{Question 2}
\textbf{Paper:} PEA \quad \textbf{Year:} 2026 \quad \textbf{Topic:} Algebra \\
\textbf{Subtopic/Concept:} Symmetric expressions in $x$ and $1/x$ \quad
\textbf{Difficulty:} Easy \quad \textbf{Status:} Verified

\subsection*{Question}
If $x^2 - 3x + 1 = 0$, then the value of $\left(x^4 + \dfrac{1}{x^4}\right)$ is
\begin{enumerate}[label=(\Alph*)]
\item 43 \item 45 \item 47 \item 49
\end{enumerate}

\subsection*{Solution}
From $x^2 - 3x + 1 = 0$ (and $x \ne 0$, since the constant term is non-zero), divide by $x$:
\[
x + \frac{1}{x} = 3.
\]
Square:
\[
x^2 + \frac{1}{x^2} = \left(x + \tfrac{1}{x}\right)^2 - 2 = 9 - 2 = 7.
\]
Square again:
\[
x^4 + \frac{1}{x^4} = \left(x^2 + \tfrac{1}{x^2}\right)^2 - 2 = 49 - 2 = 47.
\]

\subsection*{Final Answer}
\[
\boxed{\text{Option C: } 47}
\]

\subsection*{Common Trap / Note}
No major trap beyond standard calculation care.

%--------------------------------------------------
\section*{Question 3}
\textbf{Paper:} PEA \quad \textbf{Year:} 2026 \quad \textbf{Topic:} Calculus \\
\textbf{Subtopic/Concept:} Definite integrals, substitution \quad
\textbf{Difficulty:} Easy \quad \textbf{Status:} Verified

\subsection*{Question}
Let $f:\mathbb{R}\to\mathbb{R}$ satisfy $f(x^3+1) = x^6 + 4x^3 + 8$ for all
$x \in \mathbb{R}$. Then $\displaystyle \int_1^2 f(x)\,dx$ equals
\begin{enumerate}[label=(\Alph*)]
\item $\dfrac{29}{3}$ \item $10$ \item $11$ \item $\dfrac{31}{3}$
\end{enumerate}

\subsection*{Solution}
Set $u = x^3 + 1$, so $x^3 = u - 1$ and
\[
f(u) = (u-1)^2 + 4(u-1) + 8 = u^2 - 2u + 1 + 4u - 4 + 8 = u^2 + 2u + 5.
\]
Therefore
\[
\int_1^2 f(x)\,dx = \int_1^2 (x^2 + 2x + 5)\,dx
= \left[\frac{x^3}{3} + x^2 + 5x\right]_1^2
= \left(\tfrac{8}{3} + 4 + 10\right) - \left(\tfrac{1}{3} + 1 + 5\right)
= \tfrac{7}{3} + 8 = \tfrac{31}{3}.
\]

\subsection*{Final Answer}
\[
\boxed{\text{Option D: } \tfrac{31}{3}}
\]

\subsection*{Common Trap / Note}
Setting $u = x^3 + 1$ is bijective on $\mathbb{R}$, so $f$ is uniquely determined.

%--------------------------------------------------
\section*{Question 4}
\textbf{Paper:} PEA \quad \textbf{Year:} 2026 \quad \textbf{Topic:} Calculus \\
\textbf{Subtopic/Concept:} Integral of an absolute value \quad
\textbf{Difficulty:} Easy \quad \textbf{Status:} Verified

\subsection*{Question}
Determine the value of
\[
\int_0^2 |x^2 - 7x + 10|\,dx.
\]
\begin{enumerate}[label=(\Alph*)]
\item 8 \item 9 \item $\dfrac{26}{3}$ \item $-\dfrac{23}{3}$
\end{enumerate}

\subsection*{Solution}
Factor: $x^2 - 7x + 10 = (x-2)(x-5)$. On $[0,2]$, both factors satisfy
$x-2 \le 0$ and $x-5 < 0$, so their product is $\ge 0$ on $[0,2]$ (with equality
only at $x=2$). Hence $|x^2 - 7x + 10| = x^2 - 7x + 10$ on $[0,2]$, and
\[
\int_0^2 (x^2 - 7x + 10)\,dx
= \left[\frac{x^3}{3} - \frac{7x^2}{2} + 10x\right]_0^2
= \frac{8}{3} - 14 + 20 = \frac{8}{3} + 6 = \frac{26}{3}.
\]

\subsection*{Final Answer}
\[
\boxed{\text{Option C: } \tfrac{26}{3}}
\]

\subsection*{Common Trap / Note}
Verifying the sign of the quadratic on the interval is essential before dropping the absolute value.

%--------------------------------------------------
\section*{Question 5}
\textbf{Paper:} PEA \quad \textbf{Year:} 2026 \quad \textbf{Topic:} Real Analysis \\
\textbf{Subtopic/Concept:} Rolle's theorem applied to an antiderivative \quad
\textbf{Difficulty:} Moderate \quad \textbf{Status:} Verified

\subsection*{Question}
Let $f:[0,1]\to\mathbb{R}$ be continuous with $\displaystyle\int_0^1 f(t)\,dt = 1$, and let
\[
P(x) = \sum_{k=1}^n a_k x^k, \qquad \sum_{k=1}^n a_k = 1.
\]
Then there exists $c \in (0,1)$ such that
\begin{enumerate}[label=(\Alph*)]
\item $f(c) = 2\sum_{k=1}^n a_k$
\item $f(c) = \int_0^1 P(t)\,dt$
\item $f(c) = P'(c)$
\item $f(c) = P(c)$
\end{enumerate}

\subsection*{Solution}
Define
\[
\Phi(x) \;=\; \int_0^x f(t)\,dt \;-\; P(x), \qquad x \in [0,1].
\]
Then $\Phi$ is continuous on $[0,1]$, differentiable on $(0,1)$, and
\[
\Phi(0) = 0 - P(0) = 0, \qquad
\Phi(1) = \int_0^1 f(t)\,dt - P(1) = 1 - \sum_{k=1}^n a_k = 1 - 1 = 0.
\]
By Rolle's theorem, there exists $c \in (0,1)$ with $\Phi'(c) = 0$, i.e.
\[
f(c) - P'(c) = 0 \quad\Longleftrightarrow\quad f(c) = P'(c).
\]

\subsection*{Final Answer}
\[
\boxed{\text{Option C: } f(c) = P'(c)}
\]

\subsection*{Common Trap / Note}
Note $P(0)=0$ because the sum starts at $k=1$. This is essential for $\Phi(0)=0$.

%--------------------------------------------------
\section*{Question 6}
\textbf{Paper:} PEA \quad \textbf{Year:} 2026 \quad \textbf{Topic:} Real Analysis \\
\textbf{Subtopic/Concept:} Cauchy mean value theorem \quad
\textbf{Difficulty:} Moderate \quad \textbf{Status:} Verified

\subsection*{Question}
Let $f:[0,\pi/4]\to\mathbb{R}$ be continuous. Then there exists $c\in(0,\pi/4)$ such that
\begin{enumerate}[label=(\Alph*)]
\item $f(c) = \int_0^{\pi/4} f(t)\,dt$
\item $f(c) = 2\cos(2c)\,\int_0^{\pi/4} f(t)\,dt$
\item $f(c) = \sin(2c)$
\item $f(c) = \sin(2c)\,\int_0^{\pi/4} f(t)\,dt$
\end{enumerate}

\subsection*{Solution}
Define
\[
F(x) \;=\; \int_0^x f(t)\,dt, \qquad G(x) \;=\; \sin(2x), \qquad x \in [0,\pi/4].
\]
Both are continuous on $[0,\pi/4]$ and differentiable on $(0,\pi/4)$ with
$F'(x) = f(x)$ and $G'(x) = 2\cos(2x)$. Moreover $G'(x) = 2\cos(2x) \neq 0$ on
$(0,\pi/4)$, since $0 < 2x < \pi/2$ there.

By Cauchy's mean value theorem, there exists $c \in (0,\pi/4)$ such that
\[
\frac{F(\pi/4) - F(0)}{G(\pi/4) - G(0)} \;=\; \frac{F'(c)}{G'(c)},
\]
i.e.
\[
\frac{\int_0^{\pi/4} f(t)\,dt}{\sin(\pi/2) - 0} \;=\; \frac{f(c)}{2\cos(2c)},
\]
which simplifies to
\[
f(c) \;=\; 2\cos(2c)\,\int_0^{\pi/4} f(t)\,dt.
\]

\subsection*{Final Answer}
\[
\boxed{\text{Option B: } f(c) = 2\cos(2c)\int_0^{\pi/4} f(t)\,dt}
\]

\subsection*{Common Trap / Note}
The integral on the right is a constant; the $c$-dependence enters only through $\cos(2c)$.

%--------------------------------------------------
\section*{Question 7}
\textbf{Paper:} PEA \quad \textbf{Year:} 2026 \quad \textbf{Topic:} Real Analysis \\
\textbf{Subtopic/Concept:} FTC applied to an identity \quad
\textbf{Difficulty:} Easy \quad \textbf{Status:} Verified

\subsection*{Question}
Suppose $f:[a,b]\to\mathbb{R}$ is continuous and
\[
\int_a^x f(t)\,dt \;=\; \int_x^b f(t)\,dt \quad \text{for all } x \in [a,b].
\]
Then $f(x)$ must be
\begin{enumerate}[label=(\Alph*)]
\item A non-zero constant
\item A linear function
\item Identically zero
\item An odd function about $\dfrac{a+b}{2}$
\end{enumerate}

\subsection*{Solution}
Differentiate both sides with respect to $x$ using the fundamental theorem of calculus:
\[
\frac{d}{dx}\int_a^x f(t)\,dt = f(x), \qquad
\frac{d}{dx}\int_x^b f(t)\,dt = -f(x).
\]
Hence $f(x) = -f(x)$, so $2f(x) = 0$, i.e.\ $f(x) = 0$ for every $x \in [a,b]$.

\subsection*{Final Answer}
\[
\boxed{\text{Option C: identically zero}}
\]

\subsection*{Common Trap / Note}
An odd function about $\tfrac{a+b}{2}$ would satisfy the identity only at the
single point $x = \tfrac{a+b}{2}$, not for all $x$.

%--------------------------------------------------
\section*{Question 8}
\textbf{Paper:} PEA \quad \textbf{Year:} 2026 \quad \textbf{Topic:} Calculus \\
\textbf{Subtopic/Concept:} L'H\^opital's rule, symmetry of odd functions \quad
\textbf{Difficulty:} Hard \quad \textbf{Status:} Verified

\subsection*{Question}
Let $f:\mathbb{R}\to\mathbb{R}$ be a continuous odd function, vanishing exactly at one point and satisfying $f(1) = \tfrac{1}{2}$. If
\[
\lim_{x\to 1} \frac{F(x)}{G(x)} = \frac{1}{14},
\]
where
\[
F(x) = \int_{-1}^{x} f(t)\,dt, \qquad G(x) = \int_{-1}^{x} t\,|f(f(t))|\,dt,
\]
then the value of $f\!\left(\tfrac{1}{2}\right)$ is
\begin{enumerate}[label=(\Alph*)]
\item $\dfrac{1}{7}$ \item $7$ \item $-7$ \item $\dfrac{7}{2}$
\end{enumerate}

\subsection*{Solution}
\textbf{Both numerator and denominator vanish at $x=1$.}
Since $f$ is odd, $\int_{-1}^{1} f(t)\,dt = 0$, so $F(1) = 0$.
For $G(1)$, note that $f$ odd implies $f(f(-t)) = f(-f(t)) = -f(f(t))$, so
$|f(f(\cdot))|$ is even. Hence $t\,|f(f(t))|$ is an odd function of $t$, and
\[
G(1) = \int_{-1}^{1} t\,|f(f(t))|\,dt = 0.
\]
Thus the limit is of indeterminate form $0/0$, and by L'H\^opital's rule:
\[
\frac{1}{14} = \lim_{x\to 1} \frac{F'(x)}{G'(x)}
= \lim_{x\to 1} \frac{f(x)}{x\,|f(f(x))|}
= \frac{f(1)}{1\cdot |f(f(1))|}
= \frac{1/2}{|f(1/2)|}.
\]
Therefore $|f(1/2)| = 14 \cdot \tfrac{1}{2} = 7$.

\textbf{Determining the sign.}
Since $f$ is continuous and vanishes only at one point, and $f$ odd implies
$f(0)=0$, that single zero must be at $0$. Thus $f$ has constant sign on
$(0,\infty)$. From $f(1) = \tfrac{1}{2} > 0$, we conclude $f > 0$ on $(0,\infty)$,
so $f(1/2) > 0$, giving $f(1/2) = 7$.

\subsection*{Final Answer}
\[
\boxed{\text{Option B: } 7}
\]

\subsection*{Common Trap / Note}
The symmetry argument is the key step: without it the denominator $G(1)$ does not
obviously vanish and L'H\^opital cannot be applied.

%--------------------------------------------------
\section*{Question 9}
\textbf{Paper:} PEA \quad \textbf{Year:} 2026 \quad \textbf{Topic:} Statistics \\
\textbf{Subtopic/Concept:} Interpretation of zero correlation \quad
\textbf{Difficulty:} Easy \quad \textbf{Status:} Verified

\subsection*{Question}
If the correlation between $X$ and $Y$ is zero, then which of the following is always true?
\begin{enumerate}[label=(\Alph*)]
\item They are independent
\item There is no linear relationship
\item They cannot be functionally related
\item All of the above
\end{enumerate}

\subsection*{Solution}
Zero correlation is by definition the absence of any linear association.
Counterexamples eliminate the other options:
\begin{itemize}
\item Zero correlation does not imply independence (e.g.\ $Y = X^2$ with $X$ symmetric about $0$ yields zero correlation but $X, Y$ are dependent).
\item In the same example, $X$ and $Y$ are functionally related.
\end{itemize}
Thus only ``no linear relationship'' is always true.

\subsection*{Final Answer}
\[
\boxed{\text{Option B: no linear relationship}}
\]

\subsection*{Common Trap / Note}
``Uncorrelated'' is strictly weaker than ``independent'' except for jointly Gaussian variables.

%--------------------------------------------------
\section*{Question 10}
\textbf{Paper:} PEA \quad \textbf{Year:} 2026 \quad \textbf{Topic:} Regression \\
\textbf{Subtopic/Concept:} OLS passes through the mean point \quad
\textbf{Difficulty:} Easy \quad \textbf{Status:} Verified

\subsection*{Question}
In simple linear regression $Y = a + bX$, the least squares line based on $n$ data points always passes through
\begin{enumerate}[label=(\Alph*)]
\item $(0,0)$
\item $\left(\tfrac{1}{n}, \tfrac{1}{n}\right)$
\item $(1,1)$
\item None of the previous
\end{enumerate}

\subsection*{Solution}
The OLS line always passes through the point $(\bar{X}, \bar{Y})$, the sample means
of $X$ and $Y$. None of the offered points coincides with $(\bar{X}, \bar{Y})$ in general.

\subsection*{Final Answer}
\[
\boxed{\text{Option D: None of the previous}}
\]

\subsection*{Common Trap / Note}
No major trap beyond standard calculation care.

%--------------------------------------------------
\section*{Question 11}
\textbf{Paper:} PEA \quad \textbf{Year:} 2026 \quad \textbf{Topic:} Hypothesis Testing \\
\textbf{Subtopic/Concept:} Trade-off between $\alpha$ and power \quad
\textbf{Difficulty:} Easy \quad \textbf{Status:} Verified

\subsection*{Question}
Given sample size $n$, reducing the significance level $\alpha$
\begin{enumerate}[label=(\Alph*)]
\item Increases Type I error
\item Decreases Type II error
\item Increases power
\item Decreases power
\end{enumerate}

\subsection*{Solution}
By definition $\alpha = P(\text{Type I})$, so smaller $\alpha$ means fewer Type I
errors. For a fixed $n$, the rejection region shrinks, increasing the probability
of failing to reject the null when it is false, i.e.\ Type II error $\beta$ rises
and power $1-\beta$ falls.

\subsection*{Final Answer}
\[
\boxed{\text{Option D: decreases power}}
\]

\subsection*{Common Trap / Note}
Power can only be raised (for fixed effect size) by increasing $n$ or $\alpha$.

%--------------------------------------------------
\section*{Question 12}
\textbf{Paper:} PEA \quad \textbf{Year:} 2026 \quad \textbf{Topic:} Random Variables \\
\textbf{Subtopic/Concept:} Moments of a Bernoulli variable \quad
\textbf{Difficulty:} Easy \quad \textbf{Status:} Verified

\subsection*{Question}
If $X \sim \text{Bernoulli}(p)$, then $E(X^2)$ equals
\begin{enumerate}[label=(\Alph*)]
\item $p^2$ \item $p$ \item $p(1-p)$ \item $p + p^2$
\end{enumerate}

\subsection*{Solution}
Since $X \in \{0,1\}$, we have $X^2 = X$, hence $E(X^2) = E(X) = p$.

\subsection*{Final Answer}
\[
\boxed{\text{Option B: } p}
\]

\subsection*{Common Trap / Note}
Note $\operatorname{Var}(X) = p(1-p) = E(X^2) - (E X)^2$.

%--------------------------------------------------
\section*{Question 13}
\textbf{Paper:} PEA \quad \textbf{Year:} 2026 \quad \textbf{Topic:} Statistics \\
\textbf{Subtopic/Concept:} Unbiased estimation for Poisson \quad
\textbf{Difficulty:} Moderate \quad \textbf{Status:} Verified

\subsection*{Question}
For $X_1, \dots, X_n \sim \text{Poisson}(\lambda)$ (i.i.d.), which is unbiased for $\lambda^2$?
\begin{enumerate}[label=(\Alph*)]
\item $\bar{X}^2$
\item $\bar{X}$
\item $\bar{X}^2 - \dfrac{\bar{X}}{n}$
\item $\bar{X}^2 - \bar{X}$
\end{enumerate}

\subsection*{Solution}
For Poisson, $E(\bar{X}) = \lambda$ and $\operatorname{Var}(\bar{X}) = \lambda/n$. Therefore
\[
E(\bar{X}^2) = \operatorname{Var}(\bar{X}) + (E\bar{X})^2 = \frac{\lambda}{n} + \lambda^2.
\]
Consequently
\[
E\!\left(\bar{X}^2 - \frac{\bar{X}}{n}\right)
= \frac{\lambda}{n} + \lambda^2 - \frac{\lambda}{n}
= \lambda^2.
\]
Hence $\bar{X}^2 - \bar{X}/n$ is unbiased for $\lambda^2$.

\subsection*{Final Answer}
\[
\boxed{\text{Option D: } \bar{X}^2 - \tfrac{\bar{X}}{n}}
\]

\subsection*{Common Trap / Note}
Option (D)-correction subtracts $\bar X/n$, \emph{not} $\bar X$. The bias of $\bar X^2$ is $\lambda/n$, not $\lambda$.

%--------------------------------------------------
\section*{Question 14}
\textbf{Paper:} PEA \quad \textbf{Year:} 2026 \quad \textbf{Topic:} Descriptive Statistics \\
\textbf{Subtopic/Concept:} Robustness of median vs.\ mean \quad
\textbf{Difficulty:} Easy \quad \textbf{Status:} Verified

\subsection*{Question}
A dataset consists of $9$ observations arranged in increasing order
$x_1 \le x_2 \le \cdots \le x_9$ with mean $20$ and median $18$. If $x_9$ is
replaced by a very large number, which must be true?
\begin{enumerate}[label=(\Alph*)]
\item Both the mean and the median increase
\item The mean increases but the median remains unchanged
\item The median increases but the mean remains unchanged
\item Both the mean and the median remain unchanged
\end{enumerate}

\subsection*{Solution}
The median of $9$ ordered observations is $x_5$, which is unaffected by changing
the largest observation $x_9$ (the ordering is preserved). The mean equals the
sum divided by $9$, so increasing $x_9$ raises the sum and hence the mean.

\subsection*{Final Answer}
\[
\boxed{\text{Option B}}
\]

\subsection*{Common Trap / Note}
The median is robust to extreme values; the mean is not.

%--------------------------------------------------
\section*{Question 15}
\textbf{Paper:} PEA \quad \textbf{Year:} 2026 \quad \textbf{Topic:} Probability \\
\textbf{Subtopic/Concept:} Bayes' rule \quad
\textbf{Difficulty:} Easy \quad \textbf{Status:} Verified

\subsection*{Question}
A restaurant has $60\%$ vegetarian customers (of whom $30\%$ order dessert) and
$40\%$ non-vegetarian (of whom $50\%$ order dessert). Given that a randomly
selected customer ordered dessert, what is the probability they ordered vegetarian?
\begin{enumerate}[label=(\Alph*)]
\item $0.375$ \item $0.474$ \item $0.600$ \item $0.750$
\end{enumerate}

\subsection*{Solution}
By Bayes' theorem
\[
P(V \mid D) = \frac{P(D \mid V)\,P(V)}{P(D \mid V) P(V) + P(D \mid V^c) P(V^c)}
= \frac{0.3 \cdot 0.6}{0.3 \cdot 0.6 + 0.5 \cdot 0.4}
= \frac{0.18}{0.38}
= \frac{9}{19} \approx 0.474.
\]

\subsection*{Final Answer}
\[
\boxed{\text{Option B: } 0.474}
\]

\subsection*{Common Trap / Note}
No major trap beyond standard calculation care.

%--------------------------------------------------
\section*{Question 16}
\textbf{Paper:} PEA \quad \textbf{Year:} 2026 \quad \textbf{Topic:} Consumer Theory \\
\textbf{Subtopic/Concept:} Properties of preferences \quad
\textbf{Difficulty:} Moderate \quad \textbf{Status:} Verified

\subsection*{Question}
Suppose preferences are represented by $u(x,y) = x - \ln y$ for $x,y > 0$. Then the underlying preference relation must be
\begin{enumerate}[label=(\Alph*)]
\item Incomplete
\item Complete but intransitive
\item Complete and transitive but discontinuous
\item Complete, transitive and continuous but not convex
\end{enumerate}

\subsection*{Solution}
Since $u$ is a real-valued function on $\mathbb{R}_{++}^2$, the induced
preference is automatically \emph{complete} and \emph{transitive}. Since $u$ is
continuous, the preference is \emph{continuous}.

\emph{Convexity.} The upper contour set is
\[
\{(x,y): x - \ln y \ge c\} = \{(x,y): x \ge c + \ln y\}.
\]
The boundary $x = c + \ln y$ is \emph{concave} in $y$ (second derivative
$-1/y^2 < 0$), so the region lying \emph{above} it is not a convex set.
Equivalently, the Hessian
\[
\nabla^2 u = \begin{pmatrix} 0 & 0 \\ 0 & 1/y^2 \end{pmatrix}
\]
is positive semidefinite, not negative semidefinite, so $u$ is convex (not concave) in $y$ and the preference fails to be convex.

Hence the preference is complete, transitive, continuous, but not convex.

\subsection*{Final Answer}
\[
\boxed{\text{Option D}}
\]

\subsection*{Common Trap / Note}
Note that $y$ enters the utility with a \emph{negative} sign through $-\ln y$; so the good $y$ is a ``bad'' here, and indifference curves slope the wrong way for convexity.

%--------------------------------------------------
\section*{Question 17}
\textbf{Paper:} PEA \quad \textbf{Year:} 2026 \quad \textbf{Topic:} General Equilibrium \\
\textbf{Subtopic/Concept:} Exchange economy, gains from trade \quad
\textbf{Difficulty:} Moderate \quad \textbf{Status:} Draft

\subsection*{Question}
Suppose there are two agents $A$ and $B$ with utilities
\[
u_A = x_A + y_A - \tfrac{m}{2}(x_A - y_A)^2, \qquad u_B = x_B + y_B,
\]
where $x_i$ and $y_i$ are the amounts of goods $X$ and $Y$ consumed by agent
$i \in \{A,B\}$. The total endowments are $x$ and $y$. Initially $A$ owns the
entire endowment of $Y$ and $B$ owns the entire endowment of $X$. Then $A$ will
voluntarily transfer a positive amount of $Y$ to $B$ if
\begin{enumerate}[label=(\Alph*)]
\item $m > 0$ \item $m < 0$ \item $y > 2$ \item $my > \dfrac{1}{2}$
\end{enumerate}

\subsection*{Solution}
Compute $A$'s marginal rate of substitution at the endowment $(x_A, y_A) = (0, y)$:
\[
\frac{\partial u_A}{\partial x_A} = 1 - m(x_A - y_A), \qquad
\frac{\partial u_A}{\partial y_A} = 1 + m(x_A - y_A).
\]
At $(0,y)$:
\[
\text{MRS}^A_{x,y}
= \frac{\partial u_A/\partial x_A}{\partial u_A/\partial y_A}
= \frac{1 + my}{1 - my}.
\]
Agent $B$'s utility is linear with $\text{MRS}^B_{x,y} = 1$.

For $A$ to voluntarily transfer $Y$ to $B$ in exchange for $X$ (so that both
agents are made strictly better off), $A$ must value $X$ relative to $Y$ more
than $B$ does at the endowment, i.e.
\[
\text{MRS}^A_{x,y} > \text{MRS}^B_{x,y} \;=\; 1.
\]
Assuming $1 - my > 0$ (so that marginal utilities are positive at the endowment),
\[
\frac{1+my}{1-my} > 1 \quad\Longleftrightarrow\quad 1 + my > 1 - my
\quad\Longleftrightarrow\quad my > 0.
\]
Since the endowment $y > 0$, this reduces to $m > 0$.

\subsection*{Final Answer}
\[
\boxed{\text{Option A: } m > 0}
\]

\subsection*{Common Trap / Note}
The condition is on the \emph{sign} of $m$, not on its magnitude relative to $y$,
because the relevant comparison is between $A$'s MRS at the endowment and $B$'s
constant MRS of unity.

%--------------------------------------------------
\section*{Question 18}
\textbf{Paper:} PEA \quad \textbf{Year:} 2026 \quad \textbf{Topic:} Public Goods \\
\textbf{Subtopic/Concept:} Samuelson condition \quad
\textbf{Difficulty:} Moderate \quad \textbf{Status:} Needs Review

\subsection*{Question}
Two agents $A$ and $B$ have utilities
\[
u_A = x_A + y, \qquad u_B = 2x_B + y,
\]
where $x_A, x_B$ are amounts of a private good and $y$ is the amount of a public
good. Production satisfies $y + m(x_A + x_B) = 1$, with $m \in (0,1)$. Suppose a
social planner decides that one unit should be produced. For what value of $m$
will the resulting allocation be Pareto optimal?
\begin{enumerate}[label=(\Alph*)]
\item $\dfrac{4}{5}$ \item $\dfrac{4}{7}$ \item $\dfrac{4}{9}$ \item $\dfrac{4}{11}$
\end{enumerate}

\subsection*{Solution}
Marginal rates of substitution between the public good $y$ and the private good are
\[
\text{MRS}^A_{y,x} = \frac{\partial u_A/\partial y}{\partial u_A/\partial x_A} = 1,
\qquad
\text{MRS}^B_{y,x} = \frac{\partial u_B/\partial y}{\partial u_B/\partial x_B} = \tfrac{1}{2}.
\]
The marginal rate of transformation between $y$ and the private good, from
$y + m(x_A + x_B) = 1$, is $\text{MRT}_{y,x} = \tfrac{1}{m}$ (units of private
good required per additional unit of public good).

The Samuelson condition for Pareto optimality with a public good is
\[
\text{MRS}^A_{y,x} + \text{MRS}^B_{y,x} \;=\; \text{MRT}_{y,x},
\]
which yields
\[
1 + \tfrac{1}{2} = \tfrac{1}{m} \quad\Longrightarrow\quad m = \tfrac{2}{3}.
\]
This value is not among the options listed. Given the printed answer choices
$\{4/5, 4/7, 4/9, 4/11\}$, the question as transcribed is internally inconsistent
with the standard Samuelson criterion, so the problem statement (in particular
the utility specifications or the production technology) is flagged for review.

\subsection*{Final Answer}
\[
\boxed{\text{Needs Review (standard Samuelson computation yields } m = \tfrac{2}{3})}
\]

\subsection*{Common Trap / Note}
The transcribed utilities yield $m = 2/3$, which does not appear among the
options. Until the original problem statement (utility functions and/or
production technology) is verified, the correct option cannot be uniquely
determined.

%--------------------------------------------------
\section*{Question 19}
\textbf{Paper:} PEA \quad \textbf{Year:} 2026 \quad \textbf{Topic:} Market Structure \\
\textbf{Subtopic/Concept:} Monopoly with discrete buyers \quad
\textbf{Difficulty:} Moderate \quad \textbf{Status:} Verified

\subsection*{Question}
Three consumers with valuations $v_A, v_B, v_C$ each wish to buy one unit. The
monopolist has zero marginal cost and knows the valuations but must charge a
uniform price. It is given that
\[
\frac{7 v_B}{6} > v_A > v_B > \frac{8 v_C}{5} > 0.
\]
The profit-maximising price is
\begin{enumerate}[label=(\Alph*)]
\item $v_A$ \item $v_B$ \item $v_C$ \item Something in $(v_C, v_B)$
\end{enumerate}

\subsection*{Solution}
The monopolist's optimum lies at one of the buyers' valuations (any higher price
loses an additional buyer without gain, any lower price wastes surplus). Profit at each:
\begin{itemize}
\item Price $v_A$: only $A$ buys, profit $= v_A$.
\item Price $v_B$: $A$ and $B$ buy, profit $= 2 v_B$.
\item Price $v_C$: all three buy, profit $= 3 v_C$.
\end{itemize}
Compare:
\[
2 v_B \;\;\text{vs}\;\; v_A: \quad v_A < \tfrac{7 v_B}{6} < 2 v_B,
\quad\text{so}\quad 2 v_B > v_A.
\]
\[
2 v_B \;\;\text{vs}\;\; 3 v_C: \quad v_B > \tfrac{8 v_C}{5},
\quad\text{so}\quad 2 v_B > \tfrac{16}{5} v_C > 3 v_C.
\]
Thus charging $v_B$ yields the highest profit.

\subsection*{Final Answer}
\[
\boxed{\text{Option B: } p = v_B}
\]

\subsection*{Common Trap / Note}
The two inequalities $7v_B/6 > v_A$ and $5 v_B > 8 v_C$ are precisely what is
needed for $v_B$ to dominate both alternatives.

%--------------------------------------------------
\section*{Question 20}
\textbf{Paper:} PEA \quad \textbf{Year:} 2026 \quad \textbf{Topic:} Game Theory \\
\textbf{Subtopic/Concept:} Voluntary contribution to a public good, symmetric Nash equilibrium \quad
\textbf{Difficulty:} Moderate \quad \textbf{Status:} Verified

\subsection*{Question}
$n$ agents simultaneously choose private consumption $x_i$ and public-good
contribution $y_i$ to maximise
\[
u_i = x_i\, \Big(\sum_{j=1}^n y_j\Big),
\]
subject to $x_i + y_i = b$, $x_i, y_i \ge 0$. As $n$ increases, the total Nash
equilibrium expenditure on the public good
\begin{enumerate}[label=(\Alph*)]
\item Remains invariant
\item Monotonically increases and approaches $b$
\item Approaches $\dfrac{b}{n}$
\item Monotonically decreases and approaches $0$
\end{enumerate}

\subsection*{Solution}
Let $Y_{-i} = \sum_{j\ne i} y_j$. Agent $i$ chooses $y_i \in [0,b]$ to maximise
\[
u_i = (b - y_i)(y_i + Y_{-i}).
\]
The FOC is
\[
-(y_i + Y_{-i}) + (b - y_i) = 0 \quad\Longrightarrow\quad y_i = \frac{b - Y_{-i}}{2}.
\]
At a symmetric equilibrium $y_i = y^*$ and $Y_{-i} = (n-1) y^*$:
\[
y^* = \frac{b - (n-1) y^*}{2} \quad\Longrightarrow\quad (n+1) y^* = b
\quad\Longrightarrow\quad y^* = \frac{b}{n+1}.
\]
Therefore total public-good expenditure is
\[
Y \;=\; n y^* \;=\; \frac{n b}{n + 1}.
\]
The map $n \mapsto \dfrac{n}{n+1}$ is strictly increasing and approaches $1$ as
$n\to\infty$. Hence $Y$ monotonically increases and approaches $b$.

\subsection*{Final Answer}
\[
\boxed{\text{Option B}}
\]

\subsection*{Common Trap / Note}
Per-capita contribution $b/(n+1)$ shrinks, but it shrinks slower than $1/n$, so
the aggregate grows.

%--------------------------------------------------
\section*{Question 21}
\textbf{Paper:} PEA \quad \textbf{Year:} 2026 \quad \textbf{Topic:} Public Goods \\
\textbf{Subtopic/Concept:} Vertical summation of demand, Samuelson condition \quad
\textbf{Difficulty:} Easy \quad \textbf{Status:} Verified

\subsection*{Question}
$A$ and $B$ are neighbours; the snow plough cannot clear in front of $A$'s house
without clearing in front of $B$'s. Inverse demands are $P_A = 12 - q$,
$P_B = 8 - q$. Marginal cost is $16$. What is the efficient level of provision?
\begin{enumerate}[label=(\Alph*)]
\item $6$ \item $4$ \item $2$ \item $1$
\end{enumerate}

\subsection*{Solution}
Since the service is non-rival between $A$ and $B$, vertical summation gives
the aggregate marginal willingness to pay:
\[
P_A + P_B \;=\; (12 - q) + (8 - q) \;=\; 20 - 2q.
\]
Efficiency requires $\sum P_i = \text{MC}$:
\[
20 - 2q \;=\; 16 \quad\Longrightarrow\quad q^* \;=\; 2.
\]

\subsection*{Final Answer}
\[
\boxed{\text{Option C: } q^* = 2}
\]

\subsection*{Common Trap / Note}
For a public good, demands are summed \emph{vertically} (prices summed at each quantity), not horizontally.

%--------------------------------------------------
\section*{Question 22}
\textbf{Paper:} PEA \quad \textbf{Year:} 2026 \quad \textbf{Topic:} Microeconomics \\
\textbf{Subtopic/Concept:} Tax incidence \quad
\textbf{Difficulty:} Easy \quad \textbf{Status:} Verified

\subsection*{Question}
A lump-sum tax of $\mathrm{Rs.}\,1$ is paid by the buyer per unit of a competitive good.
The price paid by buyers rises by $80$ paise. Then
\begin{enumerate}[label=(\Alph*)]
\item Demand must be perfectly elastic
\item Demand must be perfectly inelastic
\item Demand and supply elasticities must both be positive but finite
\item Supply must be perfectly elastic
\end{enumerate}

\subsection*{Solution}
With elasticities $\varepsilon_D, \varepsilon_S$ (positive numbers), the share of
the tax borne by buyers is
\[
\frac{\varepsilon_S}{\varepsilon_S + \varepsilon_D}.
\]
We are told this share equals $0.8$.
\begin{itemize}
\item Perfectly elastic demand ($\varepsilon_D = \infty$) gives buyer share $0$. Rules out (A).
\item Perfectly inelastic demand ($\varepsilon_D = 0$) gives buyer share $1$. Rules out (B).
\item Perfectly elastic supply ($\varepsilon_S = \infty$) gives buyer share $1$. Rules out (D).
\end{itemize}
A buyer share strictly between $0$ and $1$ requires both elasticities to be
positive and finite.

\subsection*{Final Answer}
\[
\boxed{\text{Option C}}
\]

\subsection*{Common Trap / Note}
The incidence formula uses absolute elasticities; the buyer share is decreasing in $\varepsilon_D$ and increasing in $\varepsilon_S$.

%--------------------------------------------------
\section*{Question 23}
\textbf{Paper:} PEA \quad \textbf{Year:} 2026 \quad \textbf{Topic:} Consumer Theory \\
\textbf{Subtopic/Concept:} Engel aggregation \quad
\textbf{Difficulty:} Moderate \quad \textbf{Status:} Verified

\subsection*{Question}
A consumer always spends one-fourth of income on $X$, and the income elasticity of demand for $X$ is $5$. Is $Y$ inferior?
\begin{enumerate}[label=(\Alph*)]
\item Yes \item No
\item $Y$ is inferior iff $p_Y \ge 2 p_X$
\item $Y$ is inferior iff $p_Y \le p_X/2$
\end{enumerate}

\subsection*{Solution}
Let $s_X = 1/4$ and $s_Y = 3/4$ be the budget shares. Engel aggregation states
that the share-weighted sum of income elasticities equals $1$:
\[
s_X\, \eta_X + s_Y\, \eta_Y = 1.
\]
Substituting $s_X = 1/4$, $\eta_X = 5$, $s_Y = 3/4$:
\[
\tfrac{1}{4}\cdot 5 + \tfrac{3}{4}\,\eta_Y = 1 \quad\Longrightarrow\quad
\tfrac{3}{4}\,\eta_Y = -\tfrac{1}{4} \quad\Longrightarrow\quad \eta_Y = -\tfrac{1}{3} < 0.
\]
A negative income elasticity means $Y$ is inferior, irrespective of prices.

\subsection*{Final Answer}
\[
\boxed{\text{Option A: Yes}}
\]

\subsection*{Common Trap / Note}
Engel aggregation follows from differentiating the budget identity with respect to income.

%--------------------------------------------------
\section*{Question 24}
\textbf{Paper:} PEA \quad \textbf{Year:} 2026 \quad \textbf{Topic:} Consumer Theory \\
\textbf{Subtopic/Concept:} Aggregate price elasticity as quantity-weighted average \quad
\textbf{Difficulty:} Easy \quad \textbf{Status:} Verified

\subsection*{Question}
Half the total quantity is purchased by $75$ consumers each with price elasticity
$2$, and the other half by $25$ consumers each with price elasticity $3$. What is
the aggregate price elasticity for the $100$ consumers?
\begin{enumerate}[label=(\Alph*)]
\item Cannot be determined \item $2.25$ \item $2.5$ \item $2.75$
\end{enumerate}

\subsection*{Solution}
The market elasticity is the quantity-share-weighted average of group elasticities. With each group accounting for $50\%$ of the quantity:
\[
\varepsilon = \tfrac{1}{2}\cdot 2 + \tfrac{1}{2}\cdot 3 = 2.5.
\]

\subsection*{Final Answer}
\[
\boxed{\text{Option C: } 2.5}
\]

\subsection*{Common Trap / Note}
Weights are by \emph{quantity} (or revenue, given a common price), not by \emph{number of consumers}.

%--------------------------------------------------
\section*{Question 25}
\textbf{Paper:} PEA \quad \textbf{Year:} 2026 \quad \textbf{Topic:} International Trade \\
\textbf{Subtopic/Concept:} Two-country Cobb--Douglas equilibrium \quad
\textbf{Difficulty:} Moderate \quad \textbf{Status:} Verified

\subsection*{Question}
Two countries $H$ and $F$ produce single goods with outputs $Y_H, Y_F$. Take
$H$'s good as numeraire and $p$ as the relative price of $F$'s good. Consumers
in country $i$ spend a fraction $1/4$ of their expenditure on the foreign good.
$E_i$ denotes total expenditure of country $i$ measured in its own good.
Given $Y_H = 120$, $Y_F = 100$, $E_H = 80$ and market clearing with
$Y_H + p Y_F = E_H + p E_F$, find $p$.
\begin{enumerate}[label=(\Alph*)]
\item $8$ \item $6$ \item $4$ \item $2$
\end{enumerate}

\subsection*{Solution}
By the Cobb--Douglas expenditure shares, each country spends $3/4$ on its own good and $1/4$ on the foreign good.

\textbf{Market clearing for $H$'s good (measured in units of $H$'s good).}
Demand by $H$-consumers: $\tfrac{3}{4} E_H = \tfrac{3}{4}\cdot 80 = 60$.
Demand by $F$-consumers: a fraction $1/4$ of their total expenditure $p\,E_F$
expressed in $H$-units, i.e.\ $\tfrac{1}{4}\,p\,E_F$.
Total demand equals supply $Y_H$:
\[
60 + \tfrac{1}{4}\,p\,E_F \;=\; 120 \quad\Longrightarrow\quad p\,E_F = 240.
\]

\textbf{Budget identity.} Total income $=$ total expenditure (in the same units):
\[
Y_H + p\,Y_F \;=\; E_H + p\,E_F \quad\Longrightarrow\quad
120 + 100 p \;=\; 80 + 240 \;=\; 320.
\]
Thus $100 p = 200$, hence $p = 2$.

\subsection*{Final Answer}
\[
\boxed{\text{Option D: } p = 2}
\]

\subsection*{Common Trap / Note}
Keeping the units straight (own good vs.\ numeraire) is the only delicate step.

%--------------------------------------------------
\section*{Question 26}
\textbf{Paper:} PEA \quad \textbf{Year:} 2026 \quad \textbf{Topic:} Macroeconomics \\
\textbf{Subtopic/Concept:} Classical monetary equilibrium with full employment \quad
\textbf{Difficulty:} Moderate \quad \textbf{Status:} Verified

\subsection*{Question}
An economy has one good produced one-for-one from labour. The representative
consumer has $U\!\left(C, \tfrac{M}{P}\right) = \tfrac{3}{4}\ln C + \tfrac{1}{4}\ln\!\tfrac{M}{P}$,
labour endowment $50$, money endowment $\bar M = 100$, with constant money
supply $M = \bar M$. If the price is flexible (so labour is fully employed), the
equilibrium price is
\begin{enumerate}[label=(\Alph*)]
\item $2$ \item $4$ \item $6$ \item $8$
\end{enumerate}

\subsection*{Solution}
With linear technology and perfect competition $P = W$. The consumer's real
wealth (measured in units of the good) is
\[
\Omega = L + \frac{\bar M}{P},
\]
where $L$ is labour supplied. Cobb--Douglas preferences imply
\[
C = \tfrac{3}{4}\,\Omega, \qquad \tfrac{M}{P} = \tfrac{1}{4}\,\Omega.
\]
Full employment means $L = 50$ and goods-market clearing $C = Y = L = 50$:
\[
50 = \tfrac{3}{4}\!\left(50 + \tfrac{100}{P}\right) \;\Longleftrightarrow\;
\tfrac{200}{3} - 50 = \tfrac{100}{P}\cdot \tfrac{3}{4}\cdot\tfrac{4}{3}.
\]
Equivalently, the money-market condition
\[
\frac{100}{P} = \tfrac{1}{4}\!\left(50 + \tfrac{100}{P}\right)
\]
gives
\[
\frac{400}{P} - \frac{100}{P} = 50 \;\Longrightarrow\; \frac{300}{P} = 50
\;\Longrightarrow\; P = 6.
\]

\subsection*{Final Answer}
\[
\boxed{\text{Option C: } P = 6}
\]

\subsection*{Common Trap / Note}
Either the goods market or the money market suffices (Walras' law guarantees the other clears).

%--------------------------------------------------
\section*{Question 27}
\textbf{Paper:} PEA \quad \textbf{Year:} 2026 \quad \textbf{Topic:} Macroeconomics \\
\textbf{Subtopic/Concept:} Fixed-price (Keynesian) equilibrium \quad
\textbf{Difficulty:} Moderate \quad \textbf{Status:} Verified

\subsection*{Question}
With the same setup as Question 26 but with price fixed at $P = 10$,
the equilibrium output is
\begin{enumerate}[label=(\Alph*)]
\item $15$ \item $20$ \item $25$ \item $30$
\end{enumerate}

\subsection*{Solution}
With $P=10$, real money endowment is $\bar M/P = 10$. The money-demand
condition
\[
\frac{M}{P} = \tfrac{1}{4}\!\left(L + \tfrac{\bar M}{P}\right)
\]
becomes
\[
10 = \tfrac{1}{4}(L + 10) \;\Longrightarrow\; L + 10 = 40 \;\Longrightarrow\; L = 30.
\]
Since the price is above the market-clearing level $P^* = 6$ (Question 26), the
labour market is in excess supply and employment is demand-determined.
Output $Y = L = 30$.

\subsection*{Final Answer}
\[
\boxed{\text{Option D: } Y = 30}
\]

\subsection*{Common Trap / Note}
At a fixed price above $P^*$, output is pinned down by the money-market
(or equivalently, goods-market) condition, not by the labour supply.

%--------------------------------------------------
\section*{Question 28}
\textbf{Paper:} PEA \quad \textbf{Year:} 2026 \quad \textbf{Topic:} Growth Theory \\
\textbf{Subtopic/Concept:} Solow model, growth rate of per-capita output \quad
\textbf{Difficulty:} Moderate \quad \textbf{Status:} Verified

\subsection*{Question}
Solow economy with $Y = K^{1/2} L^{1/2}$, no depreciation, $n = 0.02$, savings
rate $s > 0$. If the steady-state capital--labour ratio is $k^* = 9$ and the
current ratio is $k = 1$, then the growth rate of $y = Y/L$ at the current date
is
\begin{enumerate}[label=(\Alph*)]
\item $0.02$ \item $0.04$ \item $0.03$ \item Indeterminable
\end{enumerate}

\subsection*{Solution}
Per-worker output is $y = k^{1/2}$. The Solow law of motion (with no depreciation) is
\[
\dot k = s\,k^{1/2} - n\,k.
\]
Steady state: $s\,(k^*)^{1/2} = n\,k^* \Rightarrow s = n\,(k^*)^{1/2} = 0.02 \cdot 3 = 0.06$.

Current growth rate of $k$:
\[
\frac{\dot k}{k} = \frac{s}{\sqrt k} - n = \frac{0.06}{1} - 0.02 = 0.04.
\]
Since $y = k^{1/2}$,
\[
\frac{\dot y}{y} = \tfrac{1}{2}\cdot \frac{\dot k}{k} = \tfrac{1}{2}\cdot 0.04 = 0.02.
\]

\subsection*{Final Answer}
\[
\boxed{\text{Option A: } 0.02}
\]

\subsection*{Common Trap / Note}
Don't forget the factor $1/2$ when converting from the growth rate of $k$ to that of $y = k^{1/2}$.

%--------------------------------------------------
\section*{Question 29}
\textbf{Paper:} PEA \quad \textbf{Year:} 2026 \quad \textbf{Topic:} Growth Theory \\
\textbf{Subtopic/Concept:} Comparative dynamics across Solow economies \quad
\textbf{Difficulty:} Moderate \quad \textbf{Status:} Verified

\subsection*{Question}
Two Solow economies $A, B$ share the same technology and $n = 0.02$. They have
$k_A = 3, k^*_A = 9$ and $k_B = 4, k^*_B = 16$. Then per capita output is growing
\begin{enumerate}[label=(\Alph*)]
\item Faster in $A$ than in $B$
\item Faster in $B$ than in $A$
\item At the same rate in $A$ and $B$
\item At an indeterminable relative rate
\end{enumerate}

\subsection*{Solution}
As in Question 28, $s_i = n\,(k^*_i)^{1/2}$. Thus $s_A = 0.02\cdot 3 = 0.06$ and
$s_B = 0.02\cdot 4 = 0.08$. The current growth rate of capital per worker is
$\dot k_i/k_i = s_i/\sqrt{k_i} - n$, so
\[
\frac{\dot k_A}{k_A} = \frac{0.06}{\sqrt 3} - 0.02 \approx 0.0346 - 0.02 = 0.0146,
\qquad
\frac{\dot k_B}{k_B} = \frac{0.08}{2} - 0.02 = 0.04 - 0.02 = 0.02.
\]
Halving these gives the growth rates of $y$:
\[
\frac{\dot y_A}{y_A} \approx 0.0073, \qquad \frac{\dot y_B}{y_B} = 0.01.
\]
Hence $y$ grows faster in $B$ than in $A$.

\subsection*{Final Answer}
\[
\boxed{\text{Option B}}
\]

\subsection*{Common Trap / Note}
Both economies are below their respective steady states, but $B$ has a higher
savings rate and a closer relative gap, so it grows faster.

%--------------------------------------------------
\section*{Question 30}
\textbf{Paper:} PEA \quad \textbf{Year:} 2026 \quad \textbf{Topic:} Intertemporal Choice \\
\textbf{Subtopic/Concept:} Log utility and the saving function \quad
\textbf{Difficulty:} Moderate \quad \textbf{Status:} Verified

\subsection*{Question}
A two-period consumer maximises
\[
U(C_1, C_2) = \log C_1 + \tfrac{1}{1+\rho}\log C_2, \qquad \rho > 0,
\]
subject to $C_1 + S = w$ and $C_2 = (1+r) S$, with $w, r > 0$. If $w$ rises by
$1\%$ and $r$ falls by $1\%$, then $S$
\begin{enumerate}[label=(\Alph*)]
\item Increases by more than $1\%$
\item Increases by $1\%$
\item Decreases by $1\%$
\item Decreases by more than $1\%$
\end{enumerate}

\subsection*{Solution}
Substitute $C_2 = (1+r)(w - C_1)$ into $U$ and differentiate w.r.t.\ $C_1$:
\[
\frac{1}{C_1} - \frac{1}{1+\rho}\cdot \frac{1+r}{C_2}\cdot(1+r)\cdot\frac{1}{1+r}
= \frac{1}{C_1} - \frac{1}{(1+\rho) C_2}\,(1+r)\cdot\frac{1}{1+r} \cdot (1+r).
\]
A cleaner derivation: the Euler equation $\dfrac{C_2}{C_1} = \dfrac{1+r}{1+\rho}$, combined with the lifetime budget constraint $C_1 + \dfrac{C_2}{1+r} = w$, gives
\[
C_1 + \frac{1}{1+r}\cdot \frac{1+r}{1+\rho}\,C_1 \;=\; w
\quad\Longrightarrow\quad C_1\!\left(1 + \frac{1}{1+\rho}\right) = w
\quad\Longrightarrow\quad C_1 = \frac{(1+\rho)\,w}{2+\rho}.
\]
Therefore
\[
S = w - C_1 = w\!\left(1 - \frac{1+\rho}{2+\rho}\right) = \frac{w}{2+\rho}.
\]
This is the well-known log-utility result: \emph{saving is proportional to wage
and independent of $r$} (income and substitution effects of an interest-rate
change cancel exactly under log preferences).

Since $S$ depends only on $w$ (not on $r$):
\begin{itemize}
\item A $1\%$ rise in $w$ raises $S$ by exactly $1\%$.
\item A $1\%$ fall in $r$ leaves $S$ unchanged.
\end{itemize}
Therefore $S$ rises by exactly $1\%$.

\subsection*{Final Answer}
\[
\boxed{\text{Option B: increases by } 1\%}
\]

\subsection*{Common Trap / Note}
Under log utility the income and substitution effects of a change in $r$ on
saving cancel out exactly; this is a special property of log preferences.

%--------------------------------------------------
\newpage
\section*{Review Flags}

\textbf{Question 18.} As transcribed (with $u_A = x_A + y$, $u_B = 2x_B + y$ and
production constraint $y + m(x_A + x_B) = 1$), the Samuelson condition yields
$m = 2/3$, which is not among the printed options $\{4/5,\,4/7,\,4/9,\,4/11\}$.
The problem statement (likely the utility functions or the production technology
in the underlying source) should be verified before finalising the option choice.

All other questions are flagged as either \textbf{Verified} or \textbf{Draft}
(see the Answer-Key Summary at the front of this booklet).

\bigskip
\hrule
\bigskip
\begin{center}
\textit{End of PEA Solutions --- ISI MSQE 2026, prepared for Statstrive.}
\end{center}

\end{document}
